\chapter{总结与展望}

\section{总结}

时空序列预测是智慧城市、智能交通、能源管理、气象预报和金融分析等领域中的基础性任务。近年来，基于深度学习和图神经网络的方法在预测精度上取得了明显进展，但现实时空数据普遍存在时间非平稳、空间非平稳、多尺度模式和变量或节点间分布异质性等复杂特性。现有方法通常使用单一尺度表示、共享参数或固定空间图结构，难以同时适应不同时间状态、不同变量和不同空间区域之间的差异。为此，本文围绕“时空多尺度建模和混合专家系统”这一主题展开研究，分别针对一般多元时间序列预测和具有显式空间图结构的交通预测两类任务，从时间维度和空间维度构建了递进的研究框架。

首先，针对一般多元时间序列预测中的时间非平稳、通道分布异质以及多尺度时序模式并存问题，本文提出了一种融合双向时间戳依赖建模与多尺度特征提取混合专家的预测框架。该方法采用通道独立路由和稀疏门控，将不同变量动态分配给合适的专家，从而避免共享参数模型对不同变量分布的一刀切处理；同时，通过多尺度下采样、趋势—季节分解以及自底向上和自顶向下融合，分别提取细粒度局部模式和宏观趋势信息。进一步地，该方法在频域中利用可学习马氏矩阵学习变量相关性，并通过时间戳双向交叉注意力显式建模历史与未来时间戳之间的双向依赖。实验结果表明，该方法在多个真实数据集上总体取得较优预测结果，消融和参数分析也验证了各模块的必要性。

其次，针对交通预测中的空间非平稳、静态路网结构与动态时空依赖共存问题，本文提出了一种融合自适应空间多尺度建模与静动态双混合专家的预测框架。该方法通过图池化以数据驱动方式构造多个空间尺度，并利用尺度专家和尺度自适应融合提取尺度特异性特征；同时使用稀疏激活的静态图专家和动态图专家分别建模显式拓扑结构和隐式时空交互。实验表明，该方法在多个交通预测基准上取得较优或接近最优的结果，可视化结果进一步说明模型能够根据高峰期和非高峰期动态调整空间尺度与专家选择。

从整体研究脉络来看，本文的两项工作并不是相互独立的算法，而是同一研究框架在两种数据形态上的自然延伸。第一项工作将多尺度建模和混合专家机制应用于时间维度，解决变量间没有显式图结构时的时间分布差异问题；第二项工作将该思想扩展到空间图结构，进一步处理空间尺度变化和静动态依赖协同问题。二者共同验证了时空多尺度建模和混合专家系统在处理非平稳、异质时空数据方面的有效性。

\section{展望}

尽管本文在时空多尺度建模和混合专家系统方面进行了一些探索，但仍存在多个方向值得进一步深入研究：

\begin{enumerate}
    \item 更丰富的空间语义和多尺度表示。本文主要利用传感器观测和路网邻接关系进行建模，尚未充分利用道路等级、车道数、兴趣点和用地类型等空间语义信息。未来可将这些语义先验作为节点特征或图结构先验引入模型，进一步丰富空间表示，提升对复杂路网和区域功能的适应能力。

    \item 更灵活的多尺度构造与专家路由。当前多尺度图构造仍依赖于预设的尺度因子和池化层数，专家数量和激活数量也需要通过实验选择。未来可以探索完全数据驱动的尺度数量学习、动态专家结构以及面向不同时空状态的细粒度路由策略，使模型能够根据输入自动调整建模复杂度和计算开销。

    \item 预训练模型与多模态信息的融合。近年来，大型语言模型和通用基础模型在时间序列分析中展现出一定潜力。未来可以研究如何将预训练模型中的通用序列知识迁移到时空预测任务，并融合文本、图像、天气和事件等多模态信息，增强模型在数据稀疏和分布变化条件下的泛化能力。

    \item 面向异常事件和长时预测的鲁棒建模。本文主要关注常规交通动态和标准预测窗口，对事故、极端天气等突发事件的建模较少。未来可引入事件感知的外生信号，并研究更长预测范围下的误差累积和不确定性量化，提高模型在非平稳、非正常交通状态下的可靠性和可解释性。
\end{enumerate}

综上所述，随着深度学习、图神经网络和大规模数据资源的持续发展，时空多尺度建模和混合专家系统有望在未来进一步提升对非平稳时空数据的建模能力，并在智能交通、城市计算等实际应用中发挥更大作用。本文的研究为这一方向提供了初步探索，也为后续构建更具适应性、鲁棒性和可解释性的时空预测模型提供了参考。
